%% Research design section 12, "principal failure conditions", and section 11's invariants read
%% as a scope statement.
%%
%% Registered in docs/preregistration.md before the experiments ran, which is
%% what lets this section be written in the indicative rather than the
%% apologetic:
%%   - asymmetric clusters compress to ratio 1.0 by construction
%%     (graph-toy-asym is in the corpus so the number gets reported);
%%   - the isomorphism check can cost more than the simulation it saves, and
%%     the registered response is to demote exact compression to a cache key;
%%   - path-model error can dominate the ranking differences.
%%
%% Also here, and not buried: every fixture in this repository is fictional;
%% the contention model is fluid, not packet-level, and does not model
%% contention BETWEEN candidates; and a simulation at 128 devices is not a
%% hardware accuracy result.
\section{Limitations}
\label{sec:limits}

% claim: C4
\paragraph{The compression has a precondition.}
Exact equivalence folds placements only where the cluster has repeated
structure; on the toy asymmetric fixture it folds nothing (ratio
\egoneRatioAsym). Where the isomorphism check costs more than the simulation it
saves, the registered response is to demote exact compression to a cache key.

% claim: C13
% claim: C14
% claim: C40
\paragraph{Some placements never simulate.}
On the smallest corpus \egthreeUnjudgedLab{} of \egthreeEmbeddingsLab{}
placements returned no verdict, so the correctness result there is a
statement about the placements the simulator judged; on the toy corpora, once
the missing profile was measured, every placement was judged. Neither of the two causes
identified in Section~\ref{sec:eval} becomes a
verdict by being understood: the KV exhaustion is detected and raised, so the
failure carries a reason but not a verdict. With an A5000 decoding, the
simulator judged loads at which the hardware already preempts
(Section~\ref{sec:eval-pd}).

% claim: C34
% claim: C3
\paragraph{The simulator is told a placement's bandwidth, not what it shares.}
Every registered hardware prediction used the earlier adapter's first-pair
figure; the correction is evaluated only post hoc (Section~\ref{sec:eval-hw}).
Contention between a candidate's own flows on one port or NIC remains
invisible to the evaluator, and closing that needs a simulator with a link
graph. The same first-pair read remains in the compression's demand labels,
in the fluid model's choice of a flow's path and in the ranker's cut margin;
correcting them would change which candidates the registered runs reached, so
they are recorded rather than changed. For a tensor-parallel group of more
than two ranks, the labels and the earlier evaluator therefore shared one
blind spot, and the zero mis-merge count there is not independent evidence
that the labels are sufficient.

% claim: C22
% claim: C30
% claim: C39
\paragraph{Inter-node contention on a disaggregated deployment is smaller than
the model predicts.}
Inter-node \emph{P/D} runs---prefill on one node, decode on the other, the KV
pulled across the NIC by NIXL~\cite{nixl}---deployed by the experiment's own harness, because the
deployment backend of HeteroPilot~\cite{heteropilot} has no router. At a load the
configuration sustains, a background stream on the same NIC lengthened the
interval from prefill response to first decoded token by \egfivePdEffect{}
ms, with a standard deviation of \egfivePdSd{} ms over three pairs, so not
distinguishable from zero, where the fluid model predicted
\egfivePdPredicted{} ms: a ratio of \egfivePdRatio{}, outside the registered
agreement band. Across two accelerator types the two further directions miss
it too (Section~\ref{sec:eval-pd}). The model enters the background as a
standing reservation of its duty cycle; post hoc, a competing flow under
processor sharing brings one direction of three inside the band and leaves the
other two well below it. So the reservation is one wrong assumption, not the
only one, and nothing here separates the rest.

% claim: C38
% claim: C21
% not-established
\paragraph{Deferred by scope.}
On hardware only P/D crosses accelerator types, an A40 and an RTX A5000
(Section~\ref{sec:eval-pd}); every tensor-parallel group stays within one.
That the slower wire moves the load knee itself was seen only in the excluded
pilot runs and is not claimed. A burst-pattern defect was fixed before any
reported run. NPU deployment (no backend exists), a tensor-parallel group spanning two
accelerator vendors, KV-cache format conversion between them, learned rankers
and mixed parallelism degrees within one group are outside this work; nothing
here shows they cannot be done.
